\section*{第 \ref{ch:g7:heart-and-circulation} 章——心脏与循环}

\begin{solution}{exo:g7:heart-and-circulation:1}
右\omterm{def:g7:heart-and-circulation:rooms}{心房}：接收从全身回流的\omterm{prop:g4:journey-of-food:absorb}{血液}。右
\omterm{def:g7:heart-and-circulation:rooms}{心室}：把\omterm{prop:g4:journey-of-food:absorb}{血液}泵至肺部。左\omterm{def:g7:heart-and-circulation:rooms}{心房}：接收
来自肺部的富氧\omterm{prop:g4:journey-of-food:absorb}{血液}。左\omterm{def:g7:heart-and-circulation:rooms}{心室}：把\omterm{prop:g4:journey-of-food:absorb}{血液}
泵向全身——四个腔中壁最厚的一个。
\end{solution}

\begin{solution}{exo:g7:heart-and-circulation:2}
它们是单向的“活门”，阻止\omterm{prop:g4:journey-of-food:absorb}{血液}倒流。
心音“咚—嗒”是两组\omterm{def:g7:heart-and-circulation:rooms}{瓣膜}的闭合：入口\omterm{def:g7:heart-and-circulation:rooms}{瓣膜}
在\omterm{def:g7:heart-and-circulation:rooms}{心室}收缩推动\omterm{prop:g4:journey-of-food:absorb}{血液}时关闭（咚），出口\omterm{def:g7:heart-and-circulation:rooms}{瓣膜}
在\omterm{def:g7:heart-and-circulation:rooms}{心室}舒张时关闭（嗒）。
\end{solution}

\begin{solution}{exo:g7:heart-and-circulation:3}
肺循环：从右\omterm{def:g7:heart-and-circulation:rooms}{心室}到肺部\omterm{def:g7:heart-and-circulation:vessels}{毛细血管}
（吸收\omterm{def:g7:what-is-respiration:respiration}{氧气}、排出\omterm{def:g7:what-is-respiration:respiration}{二氧化碳}），再回到左\omterm{def:g7:heart-and-circulation:rooms}{心房}。体循环：从左\omterm{def:g7:heart-and-circulation:rooms}{心室}
到全身\omterm{def:g7:heart-and-circulation:vessels}{毛细血管}
（输送物质、回收废物），
再回到右\omterm{def:g7:heart-and-circulation:rooms}{心房}。
\end{solution}

\begin{solution}{exo:g7:heart-and-circulation:4}
\omterm{def:g7:heart-and-circulation:vessels}{动脉}：管壁厚而有弹性，
承受压力，将\omterm{prop:g4:journey-of-food:absorb}{血液}从\omterm{def:g4:heart-and-blood:heart}{心脏}输送出去。\omterm{def:g7:heart-and-circulation:vessels}{毛细血管}：管壁仅一个\omterm{def:g5:first-look-at-cells:cell}{细胞}厚，
是一切物质交换发生的场所。\omterm{def:g7:heart-and-circulation:vessels}{静脉}：管腔宽，管壁较薄，
内有单向\omterm{def:g7:heart-and-circulation:rooms}{瓣膜}，
在低压下将\omterm{prop:g4:journey-of-food:absorb}{血液}送回\omterm{def:g4:heart-and-blood:heart}{心脏}。
\end{solution}

\begin{solution}{exo:g7:heart-and-circulation:5}
它驱动体循环——这是范围大得多的回路，
\omterm{prop:g4:journey-of-food:absorb}{血液}流向脚趾和脑再流回——因此它必须用力最大，
左\omterm{def:g7:heart-and-circulation:rooms}{心室}壁也因此最厚。
\end{solution}

\begin{solution}{exo:g7:heart-and-circulation:6}
在\omterm{def:g7:heart-and-circulation:vessels}{毛细血管}中进行。肺部\omterm{def:g7:heart-and-circulation:vessels}{毛细血管}：\omterm{def:g7:what-is-respiration:respiration}{氧气}进入，\omterm{def:g7:what-is-respiration:respiration}{二氧化碳}排出；
\omterm{prop:g7:digestion-and-nutrients:absorption}{绒毛}的\omterm{def:g7:heart-and-circulation:vessels}{毛细血管}：\omterm{def:g7:digestion-and-nutrients:families}{营养物质}进入；\omterm{def:g3:bones-and-muscles:muscle}{肌肉}的\omterm{def:g7:heart-and-circulation:vessels}{毛细血管}：燃料和\omterm{def:g7:what-is-respiration:respiration}{氧气}离开，
\omterm{def:g7:what-is-respiration:respiration}{二氧化碳}进入（\omterm{def:g1:the-five-senses:sense}{皮肤}则为：热量
散出）。
\end{solution}

\begin{solution}{exo:g7:heart-and-circulation:7}
右\omterm{def:g7:heart-and-circulation:rooms}{心室}——肺部\omterm{def:g7:heart-and-circulation:vessels}{毛细血管}（携带\omterm{def:g7:what-is-respiration:respiration}{氧气}）——
左\omterm{def:g7:heart-and-circulation:rooms}{心房}——左\omterm{def:g7:heart-and-circulation:rooms}{心室}——颈部\omterm{def:g7:heart-and-circulation:vessels}{动脉}——
脑部\omterm{def:g7:heart-and-circulation:vessels}{毛细血管}——颈部\omterm{def:g7:heart-and-circulation:vessels}{静脉}——右\omterm{def:g7:heart-and-circulation:rooms}{心房}。脑部
的交换需大量且稳定的\omterm{def:g7:what-is-respiration:respiration}{氧气}
和\omterm{prop:g7:organs-at-work:bill}{葡萄糖}，并排出\omterm{def:g7:what-is-respiration:respiration}{二氧化碳}——这是要求苛刻的
“指挥官”下达的长期指令。
\end{solution}

\begin{solution}{exo:g7:heart-and-circulation:8}
\omterm{prop:g4:heart-and-blood:pulse}{脉搏}是每次\omterm{def:g7:heart-and-circulation:rooms}{心室}收缩时\omterm{def:g7:heart-and-circulation:vessels}{动脉}管壁的弹性扩张——
高压遇上有弹性的厚壁。到达\omterm{def:g7:heart-and-circulation:vessels}{静脉}时，
这一搏动的冲击已在穿过\omterm{def:g7:heart-and-circulation:vessels}{毛细血管}时耗尽：
\omterm{def:g7:heart-and-circulation:vessels}{静脉}中是低而平稳的压强，流经柔软宽大的管道，
已无搏动可言。
\end{solution}

\begin{solution}{exo:g7:heart-and-circulation:9}
上行回流：从腿部回流的\omterm{def:g7:heart-and-circulation:vessels}{静脉}血必须
在低压下逆重力上升，\omterm{def:g7:heart-and-circulation:rooms}{瓣膜}保住每一次的推进。
站立不动的士兵，小腿\omterm{def:g3:bones-and-muscles:muscle}{肌肉}不挤压\omterm{def:g7:heart-and-circulation:vessels}{静脉}，
上行便停滞，\omterm{prop:g4:journey-of-food:absorb}{血液}淤积在腿部，
脑部供血下降——于是晕厥。
行军时的双脚就是\omterm{def:g7:heart-and-circulation:vessels}{静脉}的“泵”。
\end{solution}

\begin{solution}{exo:g7:heart-and-circulation:10}
\omterm{def:g7:heart-and-circulation:vessels}{动脉}来自体循环；在其中，\omterm{prop:g7:digestion-and-nutrients:absorption}{绒毛}的
\omterm{def:g7:heart-and-circulation:vessels}{毛细血管}摄取进餐后的\omterm{def:g7:digestion-and-nutrients:families}{营养物质}
（\cref{prop:g7:digestion-and-nutrients:absorption}）；
这是一条特殊的连接：其\omterm{def:g7:heart-and-circulation:vessels}{静脉}先流经
\omterm{prop:g7:digestion-and-nutrients:absorption}{肝脏}进行分类与储存，此后\omterm{prop:g4:journey-of-food:absorb}{血液}才汇入
通往右\omterm{def:g7:heart-and-circulation:rooms}{心房}的一般回流。
本图重述了\omterm{def:g4:journey-of-food:digestion}{消化}一章的内容。
\end{solution}

\begin{solution}{exo:g7:heart-and-circulation:11}
心腔中的\omterm{prop:g4:journey-of-food:absorb}{血液}急速流过；\omterm{def:g4:heart-and-blood:heart}{心脏}自身厚实的心肌
不停歇地工作，与任何\omterm{def:g3:bones-and-muscles:muscle}{肌肉}一样需要自己的\omterm{def:g7:heart-and-circulation:vessels}{毛细血管}
供血——因此细小的供血\omterm{def:g7:heart-and-circulation:vessels}{动脉}在\omterm{def:g4:heart-and-blood:heart}{心脏}表面\omterm{def:g1:plants-around-us:tree}{分支}。
一旦这些\omterm{def:g4:heart-and-blood:vessels}{血管}变得狭窄，便会卡住泵自身的燃料供应线：
唯一永不停歇的器官，恰恰是被它们的阻塞
所饿死的那个。
\end{solution}

\begin{solution}{exo:g7:heart-and-circulation:12}
人体赖以生存的一切物质交换——\omterm{def:g7:lungs-and-gas-exchange:alveoli}{肺泡}处的\omterm{def:g7:what-is-respiration:respiration}{氧气}、
\omterm{prop:g7:digestion-and-nutrients:absorption}{绒毛}处的\omterm{def:g7:digestion-and-nutrients:families}{营养物质}，以及每块\omterm{def:g3:bones-and-muscles:muscle}{肌肉}处的燃料与废物——
都只发生在\omterm{def:g7:heart-and-circulation:vessels}{毛细血管}壁之间；\omterm{def:g7:heart-and-circulation:vessels}{动脉}和\omterm{def:g7:heart-and-circulation:vessels}{静脉}
仅仅把\omterm{prop:g4:journey-of-food:absorb}{血液}送到这些管壁，再将其带走，
而\omterm{def:g4:heart-and-blood:heart}{心脏}只是让\omterm{prop:g4:journey-of-food:absorb}{血液}持续流动。
泵与“高速公路”的存在，都是为了那些发丝般细的“小巷”；
只要任何地方的\omterm{def:g7:heart-and-circulation:vessels}{毛细血管}交通中断，
该器官便会死亡，无论管道多么完好。
\end{solution}

\begin{solution}{pb:g7:heart-and-circulation:1}
\textbf{1.} \omterm{prop:g4:heart-and-blood:pulse}{脉搏}：\omterm{def:g7:heart-and-circulation:vessels}{动脉}搏动的次数——
在手腕处测得的\omterm{def:g4:heart-and-blood:heart}{心脏}搏动速率。血压：\omterm{def:g7:heart-and-circulation:vessels}{动脉}在搏动时与
两次搏动之间的压强——即高速公路的负荷。
听诊器中的“咚—嗒”：听到的\omterm{def:g7:heart-and-circulation:rooms}{瓣膜}闭合声——
反映心门是否完好。

\textbf{2.} 健康患者静息时\omterm{prop:g4:heart-and-blood:pulse}{脉搏}较慢，说明\omterm{def:g7:heart-and-circulation:rooms}{心室}
每次搏动能泵出更多\omterm{prop:g4:journey-of-food:absorb}{血液}：训练有素的\omterm{def:g4:heart-and-blood:heart}{心脏}
平稳地待机，且仍有余力。

\textbf{3.} \omterm{def:g7:heart-and-circulation:rooms}{瓣膜}门：入口\omterm{def:g7:heart-and-circulation:rooms}{瓣膜}在
\omterm{def:g7:heart-and-circulation:rooms}{心室}收缩（咚）时关闭，出口\omterm{def:g7:heart-and-circulation:rooms}{瓣膜}在
\omterm{def:g7:heart-and-circulation:rooms}{心室}舒张（嗒）时关闭——四扇门两两成对地关闭，
干净利落。

\textbf{4.} 右\omterm{def:g7:heart-and-circulation:rooms}{心房}、右\omterm{def:g7:heart-and-circulation:rooms}{心室}、肺
\omterm{def:g7:heart-and-circulation:vessels}{毛细血管}、左\omterm{def:g7:heart-and-circulation:rooms}{心房}、左\omterm{def:g7:heart-and-circulation:rooms}{心室}、体
\omterm{def:g7:heart-and-circulation:vessels}{动脉}、体\omterm{def:g7:heart-and-circulation:vessels}{毛细血管}（以\omterm{def:g3:bones-and-muscles:muscle}{肌肉}为例）、\omterm{def:g7:heart-and-circulation:vessels}{静脉}——
最后回到右\omterm{def:g7:heart-and-circulation:rooms}{心房}。

\textbf{5.} 每次搏动都有一部分血流从漏口
倒流回去，因此体循环每次的输送量减少；
\omterm{def:g4:heart-and-blood:heart}{心脏}以增加搏动次数、加大做功来代偿，
耐力随之下降——杂音
即是倒流损失的输送量。

\textbf{6.} 血压升高（\omterm{def:g4:heart-and-blood:vessels}{血管}僵硬、变窄，
阻碍血流冲击），\omterm{prop:g4:heart-and-blood:pulse}{脉搏}引起的\omterm{def:g7:heart-and-circulation:vessels}{动脉}扩张也
随之改变；左\omterm{def:g7:heart-and-circulation:rooms}{心室}每次搏动都要
推动僵硬的循环系统，负担加重，
心肌增厚。

\textbf{7.} \omterm{def:g4:heart-and-blood:heart}{心脏}自身供血\omterm{def:g7:heart-and-circulation:vessels}{动脉}中的
斑块堵塞，限制了永不停歇的心肌的
供能通道——这是心肌梗死的典型途径。

\textbf{8.} \omterm{def:g7:heart-and-circulation:vessels}{静脉}的单向\omterm{def:g7:heart-and-circulation:rooms}{瓣膜}与缺失的
运动\omterm{def:g3:bones-and-muscles:muscle}{肌肉}泵有关：久站不动时，小腿
\omterm{def:g3:bones-and-muscles:muscle}{肌肉}从不收缩，\omterm{def:g7:heart-and-circulation:rooms}{瓣膜}的推动作用停滞，踝部
\omterm{prop:g4:journey-of-food:absorb}{血液}淤积。行走时有节律地挤压\omterm{def:g7:heart-and-circulation:vessels}{静脉}——
缓解靠的是机械作用。

\textbf{9.} 因为体循环会消耗\omterm{prop:g4:journey-of-food:absorb}{血液}中的
\omterm{def:g7:what-is-respiration:respiration}{氧气}，只有肺循环才能重新补充。跳过
肺部，就等于让耗尽\omterm{def:g7:what-is-respiration:respiration}{氧气}的\omterm{prop:g4:journey-of-food:absorb}{血液}空车出发
去输送——双循环的顺序是先补充再输送，
由\omterm{def:g4:heart-and-blood:heart}{心脏}的两侧共同保证。

\textbf{10.} \omterm{def:g7:heart-and-circulation:vessels}{动脉}：管壁厚、富有弹性、压强高——
输出的高速公路。\omterm{def:g7:heart-and-circulation:vessels}{毛细血管}：管壁只有一个\omterm{def:g5:first-look-at-cells:cell}{细胞}厚，
遍布全身——所有交换都在此进行的输送通道。
\omterm{def:g7:heart-and-circulation:vessels}{静脉}：管腔宽、带\omterm{def:g7:heart-and-circulation:rooms}{瓣膜}、压强低——负责
回流。

\textbf{11.} 每次问诊中的问题，都恰恰
在影响\omterm{def:g7:heart-and-circulation:vessels}{毛细血管}功能之处显现：漏口使
输送不足，斑块限制了供应
管路，淤血使回流受阻。测量值、\omterm{def:g7:heart-and-circulation:rooms}{瓣膜}
与高速公路只是手段；\omterm{def:g7:heart-and-circulation:vessels}{毛细血管}中的血流
——即物质交换——才是目的。

\textbf{12.} 每天运动——\omterm{def:g4:heart-and-blood:heart}{心脏}这台泵得以增强，
高速公路保持柔韧；均衡饮食——\omterm{def:g7:heart-and-circulation:vessels}{动脉}
管壁不被斑块堵塞；不吸烟——\omterm{def:g4:heart-and-blood:vessels}{血管}
既不变窄也不变硬，\omterm{prop:g4:journey-of-food:absorb}{血液}中的
载氧位点保持畅通。
\end{solution}
